\QHeading{第1题　口算。}
\par\noindent


\begin{oralgrid}
\OralItem{1}{6.5-3.2}{3.3}
\OralItem{2}{6.5+2.4}{8.9}
\OralItem{3}{8.1+0.8}{8.9}
\OralItem{4}{2-0.3}{1.7}
\OralItem{5}{1.5-1.3}{0.2}
\OralItem{6}{2.5+3.2}{5.7}
\OralItem{7}{3.5+2.3}{5.8}
\OralItem{8}{6.5-0.6}{5.9}
\OralItem{9}{3.85-2.3}{1.55}
\OralItem{10}{2.8+4.65}{7.45}
\OralItem{11}{23.5-2.8}{20.7}
\OralItem{12}{9.5-5.7}{3.8}\end{oralgrid}
\par

\QHeading{口算思路补充（对应上表题号）}

\Answer{上面列出得数，下面按小题给出不同计算思路。计数单位法适用于小数，整数题可比较凑整、拆分或商不变等方法。}
\MethodTitle{第1小题：$6.5-3.2$}

\MethodBlock{口算思路一：按相同计数单位计算}{先看成若干个十分之一，按整数加减，最后还原成小数。}{\begin{aligned}[t]&\phantom{={}}6.5-3.2\\
&=(65-32)\div 10\\
&=33\div 10\\
&=3.3\end{aligned}}

\MethodBlock{口算思路二：补差法}{从减数先补到一个整数，再补到被减数，两段差相加。}{\begin{aligned}[t]&\phantom{={}}6.5-3.2\\
&=(4-3.2)+(6.5-4)\\
&=0.8+2.5\\
&=3.3\end{aligned}}
\MethodTitle{第2小题：$6.5+2.4$}

\MethodBlock{口算思路一：按相同计数单位计算}{先看成若干个十分之一，按整数加减，最后还原成小数。}{\begin{aligned}[t]&\phantom{={}}6.5+2.4\\
&=(65+24)\div 10\\
&=89\div 10\\
&=8.9\end{aligned}}
\MethodTitle{第3小题：$8.1+0.8$}

\MethodBlock{口算思路一：按相同计数单位计算}{先看成若干个十分之一，按整数加减，最后还原成小数。}{\begin{aligned}[t]&\phantom{={}}8.1+0.8\\
&=(81+8)\div 10\\
&=89\div 10\\
&=8.9\end{aligned}}

\MethodBlock{口算思路二：凑整再减去多加的数}{把第二个数凑成整数，再减去多加的部分。}{\begin{aligned}[t]&\phantom{={}}8.1+0.8\\
&=8.1+1-0.2\\
&=9.1-0.2\\
&=8.9\end{aligned}}
\MethodTitle{第4小题：$2-0.3$}

\MethodBlock{口算思路一：按相同计数单位计算}{先看成若干个十分之一，按整数加减，最后还原成小数。}{\begin{aligned}[t]&\phantom{={}}2-0.3\\
&=(20-3)\div 10\\
&=17\div 10\\
&=1.7\end{aligned}}

\MethodBlock{口算思路二：补差法}{从减数先补到一个整数，再补到被减数，两段差相加。}{\begin{aligned}[t]&\phantom{={}}2-0.3\\
&=(1-0.3)+(2-1)\\
&=0.7+1\\
&=1.7\end{aligned}}
\MethodTitle{第5小题：$1.5-1.3$}

\MethodBlock{口算思路一：按相同计数单位计算}{先看成若干个十分之一，按整数加减，最后还原成小数。}{\begin{aligned}[t]&\phantom{={}}1.5-1.3\\
&=(15-13)\div 10\\
&=2\div 10\\
&=0.2\end{aligned}}
\MethodTitle{第6小题：$2.5+3.2$}

\MethodBlock{口算思路一：按相同计数单位计算}{先看成若干个十分之一，按整数加减，最后还原成小数。}{\begin{aligned}[t]&\phantom{={}}2.5+3.2\\
&=(25+32)\div 10\\
&=57\div 10\\
&=5.7\end{aligned}}
\MethodTitle{第7小题：$3.5+2.3$}

\MethodBlock{口算思路一：按相同计数单位计算}{先看成若干个十分之一，按整数加减，最后还原成小数。}{\begin{aligned}[t]&\phantom{={}}3.5+2.3\\
&=(35+23)\div 10\\
&=58\div 10\\
&=5.8\end{aligned}}
\MethodTitle{第8小题：$6.5-0.6$}

\MethodBlock{口算思路一：按相同计数单位计算}{先看成若干个十分之一，按整数加减，最后还原成小数。}{\begin{aligned}[t]&\phantom{={}}6.5-0.6\\
&=(65-6)\div 10\\
&=59\div 10\\
&=5.9\end{aligned}}

\MethodBlock{口算思路二：补差法}{从减数先补到一个整数，再补到被减数，两段差相加。}{\begin{aligned}[t]&\phantom{={}}6.5-0.6\\
&=(1-0.6)+(6.5-1)\\
&=0.4+5.5\\
&=5.9\end{aligned}}
\MethodTitle{第9小题：$3.85-2.3$}

\MethodBlock{口算思路一：按相同计数单位计算}{先看成若干个百分之一，按整数加减，最后还原成小数。}{\begin{aligned}[t]&\phantom{={}}3.85-2.3\\
&=(385-230)\div 100\\
&=155\div 100\\
&=1.55\end{aligned}}

\MethodBlock{口算思路二：补差法}{从减数先补到一个整数，再补到被减数，两段差相加。}{\begin{aligned}[t]&\phantom{={}}3.85-2.3\\
&=(3-2.3)+(3.85-3)\\
&=0.7+0.85\\
&=1.55\end{aligned}}
\MethodTitle{第10小题：$2.8+4.65$}

\MethodBlock{口算思路一：按相同计数单位计算}{先看成若干个百分之一，按整数加减，最后还原成小数。}{\begin{aligned}[t]&\phantom{={}}2.8+4.65\\
&=(280+465)\div 100\\
&=745\div 100\\
&=7.45\end{aligned}}
\MethodTitle{第11小题：$23.5-2.8$}

\MethodBlock{口算思路一：按相同计数单位计算}{先看成若干个十分之一，按整数加减，最后还原成小数。}{\begin{aligned}[t]&\phantom{={}}23.5-2.8\\
&=(235-28)\div 10\\
&=207\div 10\\
&=20.7\end{aligned}}

\MethodBlock{口算思路二：补差法}{从减数先补到一个整数，再补到被减数，两段差相加。}{\begin{aligned}[t]&\phantom{={}}23.5-2.8\\
&=(3-2.8)+(23.5-3)\\
&=0.2+20.5\\
&=20.7\end{aligned}}
\MethodTitle{第12小题：$9.5-5.7$}

\MethodBlock{口算思路一：按相同计数单位计算}{先看成若干个十分之一，按整数加减，最后还原成小数。}{\begin{aligned}[t]&\phantom{={}}9.5-5.7\\
&=(95-57)\div 10\\
&=38\div 10\\
&=3.8\end{aligned}}

\MethodBlock{口算思路二：补差法}{从减数先补到一个整数，再补到被减数，两段差相加。}{\begin{aligned}[t]&\phantom{={}}9.5-5.7\\
&=(6-5.7)+(9.5-6)\\
&=0.3+3.5\\
&=3.8\end{aligned}}


\par\noindent

\QHeading{第2题　脱式计算，能简算的要简算。}

\CalcBlock[方法一]{1}{\begin{aligned}[t]&\phantom{={}}\color{black}2.78+19+7.22\\&\color{answerred}=2.78+7.22+19\\&\color{answerred}=(2.78+7.22)+19\\&\color{answerred}=10+19\\&\color{answerred}=29\end{aligned}}
\par
\MethodBlock{方法二：按运算顺序直接计算（对照）}{本法只作结果核对；本题要求简便计算，前面的简便解法仍保留。}{\begin{aligned}[t]&\phantom{={}}2.78+19+7.22\\
&=21.78+7.22\\
&=29\end{aligned}}


\CalcBlock[方法一]{2}{\begin{aligned}[t]&\phantom{={}}\color{black}6.27-1.8-3.2\\&\color{answerred}=6.27-(1.8+3.2)\\&\color{answerred}=6.27-5\\&\color{answerred}=1.27\end{aligned}}
\MethodBlock{方法二：按运算顺序直接计算（对照）}{本法只作结果核对；本题要求简便计算，前面的简便解法仍保留。}{\begin{aligned}[t]&\phantom{={}}6.27-1.8-3.2\\
&=4.47-3.2\\
&=1.27\end{aligned}}

\CalcBlock{3}{\begin{aligned}[t]&\phantom{={}}\color{black}4.76+5.24-9\\&\color{answerred}=10-9\\&\color{answerred}=1\end{aligned}}
\CalcBlock[方法一]{4}{\begin{aligned}[t]&\phantom{={}}\color{black}0.54+6.5+3.5\\&\color{answerred}=0.54+(6.5+3.5)\\&\color{answerred}=0.54+10\\&\color{answerred}=10.54\end{aligned}}
\MethodBlock{方法二：按运算顺序直接计算（对照）}{本法只作结果核对；本题要求简便计算，前面的简便解法仍保留。}{\begin{aligned}[t]&\phantom{={}}0.54+6.5+3.5\\
&=7.04+3.5\\
&=10.54\end{aligned}}

\CalcBlock[方法一]{5}{\begin{aligned}[t]&\phantom{={}}\color{black}0.78+5.5-4.5\\&\color{answerred}=0.78+(5.5-4.5)\\&\color{answerred}=0.78+1\\&\color{answerred}=1.78\end{aligned}}
\MethodBlock{方法二：按运算顺序直接计算（对照）}{本法只作结果核对；本题要求简便计算，前面的简便解法仍保留。}{\begin{aligned}[t]&\phantom{={}}0.78+5.5-4.5\\
&=6.28-4.5\\
&=1.78\end{aligned}}

\CalcBlock[方法一]{6}{\begin{aligned}[t]&\phantom{={}}\color{black}4.2+2.8-0.8\\&\color{answerred}=4.2+(2.8-0.8)\\&\color{answerred}=4.2+2\\&\color{answerred}=6.2\end{aligned}}
\MethodBlock{方法二：先处理括号或乘除，再多减再加回}{减去凑成的整数后，把多减的部分加回来。}{\begin{aligned}[t]&\phantom{={}}4.2+2.8-0.8\\
&=7-0.8\\
&=7-1+0.2\\
&=6+0.2\\
&=6.2\end{aligned}}

\par\noindent

\QHeading{第3题　解决问题。}

\Problem{① 一筐桃连筐重25.6千克，卖出一半后，连筐重13.8千克。原来桃有多重？筐有多重？}
\par


\MethodTitle{方法一}
\Answer{桃重：{\Formula{\begin{aligned}[t]&\phantom{={}}(25.6-13.8)\times2\\&=11.8\times2\\&=23.6\text{（千克）}\end{aligned}}}\qquad 篮重：$25.6-23.6=2$（千克）\par 答：原来桃重23.6千克。\par 答：篮子重2千克。}

\MethodTitle{方法二：先求筐重，再求桃重}
\Answer{把卖出一半后连筐的重量乘2，就有全部桃的重量和两个筐的重量；减去原来连筐重，剩一个筐。}
\Answer{\textbf{筐重}\par \Formula{\begin{aligned}[t]&\phantom{={}}13.8\times 2-25.6\\
&=27.6-25.6\\
&=2\text{（千克）}\end{aligned}}}
\Answer{\textbf{桃重}\par \Formula{\begin{aligned}[t]&\phantom{={}}25.6-2\\
&=23.6\text{（千克）}\end{aligned}}}
\Answer{答：原来桃重23.6千克。\par 答：筐重2千克。}
\Problem{② 如图，把一根4米长的竹竿垂直插入水池中，竹竿入泥部分长0.4米，露出水面的部分长0.83米，水深多少米？}



\MethodTitle{方法一}

\par\noindent\makebox[\linewidth][c]{
\begin{tikzpicture}[x=1bp,y=-1bp,color=black,draw=black,text=black]
\begin{scope}\clip (421.67,674.01) rectangle (556.33,779.34);
\hdhead{67}{小数加减法——信息窗2}
\hdsection{1}{125}{口算。}
\hdsection{2}{308}{脱式计算，能简算的要简算。}
\hdsection{3}{713}{解决问题。}
\q{104}{172}{180}{12.5}{$6.5-3.2=$}
\q{104}{213}{180}{12.5}{$1.5-1.3=$}
\q{104}{254}{180}{12.5}{$3.85-2.3=$}
\q{294}{172}{180}{12.5}{$6.5+2.4=$}
\q{294}{213}{180}{12.5}{$2.5+3.2=$}
\q{294}{254}{180}{12.5}{$2.8+4.65=$}
\q{488}{172}{180}{12.5}{$8.1+0.8=$}
\q{488}{213}{180}{12.5}{$3.5+2.3=$}
\q{488}{254}{180}{12.5}{$23.5-2.8=$}
\q{674}{172}{180}{12.5}{$2-0.3=$}
\q{674}{213}{180}{12.5}{$6.5-0.6=$}
\q{674}{254}{180}{12.5}{$9.5-5.7=$}
\q{104}{352}{234}{12.5}{$2.78+19+7.22$}
\q{361}{352}{232}{12.5}{$6.27-1.8-3.2$}
\q{614}{352}{220}{12.5}{$4.76+5.24-9$}
\q{104}{524}{234}{12.5}{$0.54+6.5+3.5$}
\q{361}{524}{232}{12.5}{$0.78+5.5-4.5$}
\q{614}{524}{220}{12.5}{$4.2+2.8-0.8$}
\q{103}{760}{717}{12.5}{① 一筐桃连筐重25.6千克，卖出一半后，连筐重13.8千克。原来桃有}
\q{103}{803}{717}{12.5}{　多重？筐有多重？}


\begin{scope}[scale=.6666667]
\fill[black!65] (640,1132) rectangle (806,1145);
\draw[hdoutline] (640,1020)--(640,1145)--(806,1145)--(806,1020);
\draw[black,line width=.55pt] (640,1050)--(806,1050);
\foreach \yy in {1054,1064,...,1124} \draw[black!75,dashed,line width=.35pt] (642,\yy)--(804,\yy);
\draw[hdoutline,fill=white] (719,1024)--(721,1143)--(727,1143)--(725,1024)--cycle;
\draw[black,line width=.4pt] (712,1050)--(712,1132);
\node[font=\normalsize] at (701,1091) {?};
\draw[black,line width=.4pt] (731,1024)--(731,1050);
\node[anchor=west,font=\normalsize] at (735,1034) {0.83米};
\draw[black,line width=.4pt] (732,1132)--(738,1157);
\node[anchor=west,font=\normalsize] at (738,1157) {0.4米};
\end{scope}

\begin{scope}[color=answerred,draw=answerred,text=answerred]\ans{144.00}{112.67}{44.00}{10.5}{$3.3$}
\ans{144.00}{139.33}{44.00}{10.5}{$0.2$}
\ans{144.00}{167.33}{44.00}{10.5}{$1.55$}
\ans{280.67}{112.67}{44.00}{10.5}{$8.9$}
\ans{280.67}{139.33}{44.00}{10.5}{$5.7$}
\ans{280.67}{167.33}{44.00}{10.5}{$7.45$}
\ans{404.67}{112.67}{44.00}{10.5}{$8.9$}
\ans{404.67}{139.33}{44.00}{10.5}{$5.8$}
\ans{404.67}{167.33}{44.00}{10.5}{$20.7$}
\ans{520.00}{112.67}{44.00}{10.5}{$1.7$}
\ans{520.00}{139.33}{44.00}{10.5}{$5.9$}
\ans{520.00}{167.33}{44.00}{10.5}{$3.8$}
\ans{69.33}{257.33}{156.00}{10.5}{$\begin{aligned}[t]&=2.78+7.22+19\\&=(2.78+7.22)+19\\&=10+19\\&=29\end{aligned}$}
\ans{240.67}{257.33}{154.67}{10.5}{$\begin{aligned}[t]&=6.27-(1.8+3.2)\\&=6.27-5\\&=1.27\end{aligned}$}
\ans{409.33}{257.33}{146.67}{10.5}{$\begin{aligned}[t]&=10-9\\&=1\end{aligned}$}
\ans{69.33}{371.33}{156.00}{10.5}{$\begin{aligned}[t]&=0.54+(6.5+3.5)\\&=0.54+10\\&=10.54\end{aligned}$}
\ans{240.67}{371.33}{154.67}{10.5}{$\begin{aligned}[t]&=0.78+(5.5-4.5)\\&=0.78+1\\&=1.78\end{aligned}$}
\ans{409.33}{371.33}{146.67}{10.5}{$\begin{aligned}[t]&=4.2+(2.8-0.8)\\&=4.2+2\\&=6.2\end{aligned}$}
\ans{85.33}{550.00}{461.33}{10.5}{桃重：{\setlength{\jot}{-2pt}$\begin{aligned}[t]&\phantom{={}}(25.6-13.8)\times2\\&=11.8\times2\\&=23.6\text{（千克）}\end{aligned}$}\qquad 篮重：$25.6-23.6=2$（千克）\\答：原来桃重23.6千克。\\答：篮子重2千克。}

\draw[answerred,line width=.7pt,<->] (545,700) -- (545,754);
\ans{506.00}{713.33}{45.33}{8.5}{$2.77$米}
\end{scope}
\end{scope}\pgfresetboundingbox\path[use as bounding box] (421.67,674.01) rectangle (556.33,779.34);
\end{tikzpicture}}\par
\Answer{{\Formula{\begin{aligned}[t]&\phantom{={}}4-0.4-0.83\\&=3.6-0.83\\&=2.77\text{（米）}\end{aligned}}}\par 答：水深2.77米。}

\MethodTitle{方法二：先合计不在水里的两部分}
\Answer{入泥和露出水面的部分都不是水深，先合计后从竹竿全长中扣除。}
\Answer{\Formula{\begin{aligned}[t]&\phantom{={}}4-(0.4+0.83)\\
&=4-1.23\\
&=2.77\text{（米）}\end{aligned}}}
\Answer{答：水深2.77米。}
